\documentclass[10pt,a4paper]{amsart}
\usepackage[margin=19mm]{geometry}
\usepackage{amsmath,amssymb,mathtools}
\usepackage[scr=rsfs,cal=euler]{mathalfa}
\usepackage{xfrac}
\usepackage[unicode,hidelinks,bookmarks=false]{hyperref}
\newcommand{\ie}{i.e.\ }
\newcommand{\cf}{cf.\ }
\newcommand{\resp}{respectively\ }
\newcommand{\Subsection}{\subsection}
\providecommand{\nobreakdash}{\nobreak\hbox{-}}
\setlength{\emergencystretch}{2em}
\multlinegap0pt
\usepackage{amssymb}
\usepackage{float}
\usepackage{mathtools}
\newtheorem{proposition}{Proposition}
\newtheorem{theorem}{Theorem}
\title{A large-deviation principle for the empirical distribution of a regular branching random walk}
\author{Shuxiong Zhang, Yaping Zhu}
\date{Source: arXiv:2609.08714v1 (CC BY 4.0)}
\begin{document}
\maketitle

\noindent\emph{Source and licence.} A large-deviation principle for the empirical distribution of a regular branching random walk. Version arXiv:2609.08714v1 (CC BY 4.0), distributed by arXiv under the Creative Commons Attribution 4.0 International licence, http://creativecommons.org/licenses/by/4.0/, as recorded in the arXiv licence metadata for this exact version and shown on the abstract page https://arxiv.org/abs/2609.08714v1. Full licence notice: \texttt{licenses/arxiv-2609.08714-CC-BY-4.0.txt}.\par\smallskip

\noindent\textit{Editorial scope.} Counted unit: the article's theorem thm:dimension-reduction (source0.tex lines 2169--2193). The article's complete original proof of it is delivered with it (source0.tex lines 2194--2353, one \texttt{proof} environment of 160 lines). No mathematics is written, repaired or re-derived: the statement and the proof are the article's own lines, in the article's own notation, and no retained line is substituted or re-flowed.  Retained uncounted as the source-local context the proof needs: lines 465--470; lines 715--728; lines 1855--1862.  Nothing was omitted from any retained span, so the package carries no omission record.  No retained line contains a citation, so the wrapper carries no bibliography and no bibliography entry.  No retained line contains a figure environment, an image-inclusion command or drawing code, so no image is delivered and the package declares no asset dependency.  Editorial formatting only, in addition: an AMS \texttt{amsart} preamble that restates the article's own declarations and macros used by the retained lines, namely the environments \texttt{proposition}, \texttt{theorem} on separate counters, together with the standard packages those lines need.  The retained environments print in this document as Proposition 1, Theorem 1 in order of appearance, whereas the article prints the same items with the numbering of its own full document.  The complete native source of the article as distributed in the arXiv e-print for this exact version is bundled unchanged as \texttt{sources/209738/source0.tex}.
\begin{equation}\label{eq:QA}
  Q_A(q)
  :=
  \inf\{\mathcal{I_{\alpha}}(h):h\in\mathcal{H}_b,\ \Psi_A(h)=q\},
  \qquad q\in[0,1].
\end{equation}

\begin{align}\label{eq:path-L1}
  \int_{\partial\mathbb{T}_b}
  \sum_{k\geq1}|h_{\xi|k}|\,\mathbf m_b(\mathrm{d}\xi)
  &=
  \sum_{k\geq1}\int_{\partial\mathbb{T}_b}
  |h_{\xi|k}|\,\mathbf m_b(\mathrm{d}\xi)\cr
   &=
  \sum_{k\geq1}\sum_{|v|=k}\int_{[v]}
  |h_{v}|\,\mathbf m_b(\mathrm{d}\xi)\cr
  &=
  \sum_{k\geq1}\sum_{|v|=k}
  |h_v|\mathbf m_b([v])\cr
  &=\sum_{k\geq1}b^{-k}\sum_{|v|=k}|h_v|=\|h\|_b.
\end{align}

\begin{proposition}
\label{cor:no-suboptimal-maxima}
If every local maximizer of \(g_A\) is global, then
\[
  Q_A(q+)=Q_A(q)
  \qquad\text{for every }q\in(\nu(A),M_A).
\]
\end{proposition}

\begin{theorem}\label{thm:dimension-reduction}
Suppose that, for some \(\sigma\in\{-1,1\}\), the map
\(t\mapsto g_A(\sigma t)\) is strictly increasing on
 \(\mathbb{R}\). Fix $p\in(\nu(A),M_A)$. Then \(Q_A(p+)=Q_A(p)\), and the following assertions hold.\\
(i) If \(0<\alpha\le1\), then
\begin{equation}\label{eq:first-generation}
  Q_A(p)
  =\lambda\min\left\{
      \sum_{i=1}^b t_i^\alpha:
      t_i\ge0,\quad
      \frac1b\sum_{i=1}^b g_A(\sigma t_i)\ge p
    \right\}.
\end{equation}
The minimum is attained.\\
(ii) If $\alpha\geq1$ and \(t\mapsto g_A(\sigma t)\) is concave on
\([0,\infty)\), then
\[
  Q_A(p)=
  \begin{cases}
    \lambda b I_A(p), & \alpha=1,\\[3pt]
    \lambda\bigl(b^{1/(\alpha-1)}-1\bigr)^{\alpha-1}I_A(p)^\alpha,
      & \alpha>1.
  \end{cases}
\]
\end{theorem}

\begin{proof}
We first give a more restricted representation for $Q_A(p)$. Since \(g_A\) has no local maximizer, Proposition~\ref{cor:no-suboptimal-maxima} gives $Q_A(p+)=Q_A(p).$ The monotonicity assumption and \(g_A(0)=\nu(A)<p\) give
\[
  I_A(p)=\min\{t\ge0:g_A(\sigma t)\ge p\},
  \qquad
  g_A(\sigma I_A(p))=p.
\]
Since $Q_A$ is increasing on $[\nu(A),M_A)$, by the definition of $Q_A$ (see \eqref{eq:QA}), we can write
\begin{align}\label{45gtrh6u}
  Q_A(p)=\inf\{\mathcal{I_{\alpha}}(h):h\in\mathcal{H}_b,\ \Psi_A(h)\ge p\}.
\end{align}
Set $x+:=\max\{x,0\}$. For a finite-energy field \(h\) satisfying \(\Psi_A(h)\ge p\), for ${v\in\mathbb{T}_b^{\circ}}$, put
\[
  y_v:=\sigma h_v,
  \quad
  \widetilde h_v:=\sigma y_v^+.
\]
Since \(\lvert\widetilde h_v\rvert\le\lvert h_v\rvert\) and
\(\widetilde h\in\mathcal{H}_b\), it follows that
\begin{equation}\label{eq:energy-truncation}
  \mathcal{I_{\alpha}}(\widetilde h)
  =\lambda\sum_{v\ne\varnothing}(y_v^+)^\alpha
  \le \lambda\sum_{v\ne\varnothing}\lvert y_v\rvert^\alpha
  =\mathcal{I_{\alpha}}(h).
\end{equation}
Let $y=(y_v)_{v\in\mathbb{T}_b}$ and $y^+=(y^+_v)_{v\in\mathbb{T}_b}$ be fields with $y_{\varnothing}=y^+_{\varnothing}=0$. For $\mathbf m_b$-almost every \(\xi\), we have
\[
  H_{\widetilde h}(\xi)=\sigma H_{y^+}(\xi),
  \qquad
  H_h(\xi)=\sigma H_y(\xi),
  \qquad
  H_{y^+}(\xi)\ge H_y(\xi).
\]
Hence,
\begin{align}
  \Psi_A(\widetilde h)
  &=\int_{\partial\mathbb{T}_b}g_A\bigl(\sigma H_{y^+}(\xi)\bigr)\,\mathbf m_b(\mathrm{d}\xi)\cr
  &\ge\int_{\partial\mathbb{T}_b}g_A\bigl(\sigma H_y(\xi)\bigr)\,\mathbf m_b(\mathrm{d}\xi)=\Psi_A(h)\ge p.\nonumber
\end{align}
This, together with \eqref{eq:energy-truncation} and \eqref{45gtrh6u}, implies that we can write
\begin{equation}\label{eq:positive-y}
 Q_A(p)=\inf\left\{\mathcal{I_{\alpha}}(y):y\in\mathcal{H}_b,~y_v\geq0\ \text{for}\ {v\in\mathbb{T}_b},\ \int_{\partial\mathbb{T}_b}g_A\bigl(\sigma H_{y}(\xi)\bigr)\,\mathbf m_b(\mathrm{d}\xi)\ge p\right\}.
\end{equation}
\par
 We are ready to prove (i). Define
\begin{equation}\label{eq:Mp}
  M_p:=\lambda\inf\left\{
      \sum_{i=1}^b t_i^\alpha:
      t_i\ge0,\quad
      \frac1b\sum_{i=1}^b g_A(\sigma t_i)\ge p
    \right\}.
\end{equation}
We first prove $Q_A(p)\geq M_p$. Let \(y\) be feasible for the closed-constraint problem \eqref{eq:positive-y}.
Since \(0<\alpha\le1\), it follows that for any $i=1,...,b$,
\begin{equation}\label{eq:subadditivity}
\Big(\sum_{v\succeq i}y_v\Big)^\alpha
  \le\sum_{v\succeq i}y_v^\alpha<\infty,
\end{equation}
where $v\succeq i$ means $v=i$ or $v$ is a descendant of $i$. This, combined with  $y_v\geq0$, yields that if $\xi|1=i$, then
\begin{equation}\label{eq:ray-bound}
  H_y(\xi)=\sum_{k\ge1}y_{\xi|k}\le \sum_{v\succeq i}y_v<\infty.
\end{equation}
This, combined with the fact that $\mathbf m_b([i])=\frac{1}{b}$, concludes that
\begin{align}\label{eq:compressed-feasible}
\frac1b\sum_{i=1}^b g_A\Big(\sigma t_i\Big)
  &=\frac1b\sum_{i=1}^b g_A\Big(\sigma \sum_{v\succeq i}y_v\Big)\cr
  &\ge\sum^{b}_{i=1}\int_{[i]}g_A\bigl(\sigma H_y(\xi)\bigr)\,
       \mathbf m_b(\mathrm{d}\xi)\cr
  &=\int_{\partial\mathbb{T}_b}g_A\bigl(\sigma H_y(\xi)\bigr)\,
       \mathbf m_b(\mathrm{d}\xi)\geq p,
\end{align}
where $t_i=\sum_{v\succeq i}y_v$, $i=1,...,b$.
It follows that
\begin{equation}\label{eq:lower-bound-i}
  M_p\leq\lambda\sum_{i=1}^bt_i^{\alpha}
  =\lambda\sum_{i=1}^b(\sum_{v\succeq i}y_v)^\alpha
  \le\lambda\sum_{v\ne\varnothing}y_v^\alpha
  =\mathcal{I_{\alpha}}(y).
\end{equation}
This, together with \eqref{eq:positive-y}, yields
\(M_p\leq Q_A(p)\). It remains to prove \(Q_A(p)\leq M_p\).
The feasible set in \eqref{eq:Mp} is nonempty and closed, and its intersection with
$\{(t_1,...,t_b):\sum_{i=1}^b t_i^\alpha\leq L,t_i\ge0\}$ is compact for any $L>0$. Hence the infimum in \eqref{eq:Mp} is attained. Let \((t_1^*,\ldots,t_b^*)\) minimize
\eqref{eq:Mp}, and set
\begin{equation}\label{eq:first-gen-field}
  y_i^*:=t_i^*\quad(\lvert i\rvert=1),
  \qquad
  y_v^*:=0\quad(\lvert v\rvert\ge2).
\end{equation}
Then
\begin{equation}\label{eq:upper-bound-i}
  \int_{\partial\mathbb{T}_b}g_A\bigl(\sigma H_{y^*}(\xi)\bigr)\,
  \mathbf m_b(\mathrm{d}\xi)
  =\frac1b\sum_{i=1}^bg_A(\sigma t_i^*)\ge p,
  \quad
  \mathcal{I_{\alpha}}(y^*)=\lambda\sum_{i=1}^b(t_i^*)^\alpha=M_p.
\end{equation}
This entails that \(Q_A(p)\leq M_p\), and proves \textup{(i)}.
\par
We proceed to prove \textup{(ii)}. Let $y$ be feasible for the closed-constraint problem in \eqref{eq:positive-y}. Observe that by concavity of $g_A(\sigma\cdot)$ on $[0,\infty)$ and Jensen's inequality,
\begin{equation}\label{eq:Jensen}
\begin{aligned}
  p\le \int_{\partial\mathbb{T}_b}g_A\left(\sigma H_y(\xi)\right)\,\mathbf m_b(\mathrm{d}\xi)\leq g_A\left(\sigma\int_{\partial\mathbb{T}_b}H_y(\xi)\,\mathbf m_b(\mathrm{d}\xi)\right).
\end{aligned}
\end{equation}
By using similar arguments in \eqref{eq:path-L1}, it follows that
\begin{equation}\label{eq:weighted-lower}
  \sum_{v\ne\varnothing}b^{-\lvert v\rvert}y_v
  =\int_{\partial\mathbb{T}_b}H_y(\xi)\,\mathbf m_b(\mathrm{d}\xi)
  \ge I_A(p),
\end{equation}
where the inequality follows from \eqref{eq:Jensen} and increasing property of \(g_A(\sigma\,\cdot)\). We first deal with
\(\alpha=1\) in (ii).  Since \(b^{-\lvert v\rvert}\le b^{-1}\), by \eqref{eq:weighted-lower}, it follows that
\begin{align}\label{67i89ou71q}
  Q_A(p)&=\inf\left\{\lambda\sum_{v\ne\varnothing}y_v:
    y\in\mathcal{H}_b,~y_v\geq0\ \text{for}\ {v\in\mathbb{T}_b},\ \int_{\partial\mathbb{T}_b}g_A\bigl(\sigma H_{y}(\xi)\bigr)\,\mathbf m_b(\mathrm{d}\xi)\ge p\right\}\cr
    &\geq\lambda bI_A(p).
\end{align}
Moreover, by setting \(y_v=I_A(p)\) for $|v|=1$ and \(y_v=0\) for $|v|\neq1$, we can show that
$$Q_A(p)\leq\lambda bI_A(p).$$
This, combined with \eqref{67i89ou71q}, concludes the case \(\alpha=1\) in \textup{(ii)}.

\par
We proceed to deal with \(\alpha>1\) in (ii). By H\"older's inequality,
\begin{align}\label{eq:Holder}
  I_A(p)
  \le\sum_{v\ne\varnothing}b^{-\lvert v\rvert}y_v&\le\left(\sum_{v\ne\varnothing}
       (b^{-\lvert v\rvert})^{\alpha/(\alpha-1)}\right)^{(\alpha-1)/\alpha}
       \left(\sum_{v\ne\varnothing}y_v^\alpha\right)^{1/\alpha}\cr
  &=\bigl(b^{1/(\alpha-1)}-1\bigr)^{-(\alpha-1)/\alpha}
       \left(\sum_{v\ne\varnothing}y_v^\alpha\right)^{1/\alpha}.
\end{align}
The two inequalities in \eqref{eq:Holder} can be attained by setting
\[
  \bigl(I_A(p)(1-r)r^{-1}\bigr)^\alpha
  \bigl(b^{-\lvert v\rvert}\bigr)^{\alpha/(\alpha-1)}=y_v^\alpha,
  \qquad r=b^{-1/(\alpha-1)}.
\]
Thus, we can set
\[
  h_v^*=\sigma I_A(p)(1-r)r^{k-1},
  \quad |v|=k,~k\geq1.
\]
For every \(\xi\in\partial\mathbb{T}_b\),
\[
  H_{h^*}(\xi)
  =\sigma I_A(p)(1-r)\sum_{k\geq1}r^{k-1}
  =\sigma I_A(p),
\]
and therefore \(\Psi_A(h^*)=g_A(\sigma I_A(p))=p\). Moreover,
\begin{align*}
  \mathcal{I_{\alpha}}(h^*)
  &=\lambda I_A(p)^\alpha(1-r)^\alpha
    \sum_{k\geq1}b^k r^{\alpha(k-1)}\\
  &=\lambda b(1-r)^{\alpha-1}I_A(p)^\alpha\\
  &=\lambda\bigl(b^{1/(\alpha-1)}-1\bigr)^{\alpha-1}
    I_A(p)^\alpha.
\end{align*}
Together with \eqref{eq:Holder}, this proves the case \(\alpha>1\) in assertion \textup{(ii)}.
\end{proof}

\end{document}
